\documentclass[11pt]{article}
\usepackage[a4paper,margin=1in]{geometry}
\usepackage{amsmath,amssymb}
\usepackage[T1]{fontenc}
\usepackage{lmodern}
\usepackage{microtype}
\usepackage{xcolor}
% Footnotes inside \parbox (used below) are otherwise trapped; savenotes emits them
% at the end of the wrapper so they appear on the page foot as usual.
\usepackage{footnote}

\setlength{\parindent}{0pt}
\setlength{\parskip}{0.5em}

\newcommand{\promptbox}[1]{%
  \begin{savenotes}%
  \begingroup\setlength{\fboxsep}{10pt}\noindent
  \colorbox{blue!8}{\parbox{\dimexpr\linewidth-2\fboxsep\relax}{\textbf{Prompt. }#1}}
  \endgroup
  \end{savenotes}%
}
\newcommand{\resultbox}[1]{%
  \begin{savenotes}%
  \begingroup\setlength{\fboxsep}{10pt}\noindent
  \colorbox{purple!8}{\parbox{\dimexpr\linewidth-2\fboxsep\relax}{\textbf{Result. }#1}}
  \endgroup
  \end{savenotes}%
}
\newcommand{\examplebox}[1]{%
  \begin{savenotes}%
  \begingroup\setlength{\fboxsep}{10pt}\noindent
  \colorbox{green!8}{\parbox{\dimexpr\linewidth-2\fboxsep\relax}{\textbf{Example. }#1}}
  \endgroup
  \end{savenotes}%
}

\title{\textbf{Book 1 --- Lecture 2 --- Exercise 1}\\[2mm]
\large Derivatives of Four Introductory Functions}
\author{}
\date{}

\begin{document}
\maketitle

\promptbox{Calculate the derivatives\footnote{A \emph{derivative} measures how
sensitive a function's output is to a small change in its input; on a graph, it is
the slope of the tangent line at a point. The prime notation \(f'\) means ``the
derivative of \(f\) with respect to the variable named in context.''} of each of
these functions:}
\[
\begin{aligned}
f(t) &= t^4 + 3t^3 - 12t^2 + t - 6, \\
g(x) &= \sin(x) - \cos(x), \\
\theta(\alpha) &= e^\alpha + \alpha\ln(\alpha), \\
x(t) &= \sin^2(t) - \cos(t).
\end{aligned}
\]

\section*{\(f(t)\) (polynomial)}
Differentiate term by term using \(\dfrac{d}{dt}t^n = n t^{n-1}\):\footnote{This
is the \emph{power rule}. Sums and constant multiples of terms differentiate
term-by-term (\emph{linearity} of the derivative operator).}
\[
\begin{aligned}
f'(t) &= 4t^3 + 9t^2 - 24t + 1.
\end{aligned}
\]

\section*{\(g(x)\) (sine and cosine)}
Using \((\sin x)'=\cos x\) and \((\cos x)'=-\sin x\),\footnote{These are the basic
trigonometric derivatives with respect to the angle \(x\) (measured in radians,
as is standard in calculus and physics).}
\[
g'(x) = \cos(x) + \sin(x).
\]

\section*{\(\theta(\alpha)\) (exponential plus \(\alpha\ln\alpha\))}
Here \(\alpha>0\) so that \(\ln(\alpha)\) is defined.\footnote{\(\ln\) is the
natural logarithm (base \(e\)). It requires a positive argument \(\alpha>0\).}
We have \((e^\alpha)'=e^\alpha\),\footnote{The exponential \(e^x\) equals its own
derivative: \(\frac{d}{dx}e^x=e^x\).} and by the product rule,\footnote{The
\emph{product rule}: if \(u\) and \(v\) are differentiable functions of the same
variable, then \((uv)'=u'v+uv'\).}
\[
\frac{d}{d\alpha}\bigl(\alpha\ln(\alpha)\bigr)
= 1\cdot\ln(\alpha) + \alpha\cdot\frac{1}{\alpha}
= \ln(\alpha) + 1.
\]
Therefore
\[
\theta'(\alpha) = e^\alpha + \ln(\alpha) + 1.
\]

\section*{\(x(t)\) (chain rule on \(\sin^2 t\))}
Write \(\sin^2(t) = \bigl(\sin t\bigr)^2\). By the chain rule,\footnote{The
\emph{chain rule} composes rates of change: if \(y=h(u)\) and \(u=k(t)\), then
\(\frac{dy}{dt}=h'(u)\,k'(t)\). Here \(h(u)=u^2\) and \(u=\sin t\).}
\[
\frac{d}{dt}\sin^2(t) = 2\sin(t)\cos(t) = \sin(2t).
\]\footnote{The identity \(2\sin t\cos t=\sin(2t)\) is the sine double-angle formula.}
Also \(\dfrac{d}{dt}(-\cos t)=\sin t\).\footnote{Constants factor out of
derivatives, and \((\cos t)'=-\sin t\), so \(\frac{d}{dt}(-\cos t)=-(-\sin t)=\sin t\).}
Hence
\[
x'(t) = 2\sin(t)\cos(t) + \sin(t)
= \sin(t)\bigl(2\cos(t)+1\bigr).
\]

\resultbox{Summarizing,
\[
\begin{aligned}
f'(t) &= 4t^3 + 9t^2 - 24t + 1,\\
g'(x) &= \cos(x) + \sin(x),\\
\theta'(\alpha) &= e^\alpha + \ln(\alpha) + 1 \quad (\alpha>0),\\
x'(t) &= \sin(t)\bigl(2\cos(t)+1\bigr).
\end{aligned}
\]}

\examplebox{At \(\alpha=1\): \(\theta'(1)=e^1+\ln(1)+1=e+1\), since \(\ln(1)=0\).}

\end{document}
